\documentclass[12pt]{article}
\usepackage[utf8]{inputenc}
\usepackage{amsmath, amssymb, bm}
\usepackage{geometry}
\geometry{a4paper, margin=25mm}

\title{SC-APM-QLDPC符号の構築に向けた研究提案：\\アフィン置換行列と空間結合の統合}
\author{量子誤り訂正研究室}
\date{\today}

\begin{document}

\maketitle

\section{研究の目的}
本提案は、アフィン置換行列（APM）を用いた量子LDPC（QLDPC）符号の構成法と、空間結合（SC）構造を統合した新たな符号ファミリ「SC-APM-QLDPC符号」を提案するものである。APMによる高い最小距離の確保と、空間結合によるしきい値飽和特性を組み合わせ、次世代の量子計算機に適した符号の構築を目指す。

\section{背景理論}
近年のQLDPC符号の研究において、以下の2つのアプローチが重要である。
\begin{itemize}
    \item \textbf{APMによる直交性の打破}: 親行列の直交性制約を緩和し、アクティブな部分のみで直交性を保証するActive Orthogonalityの導入により、低重みの論理演算子を回避し、内周（girth）と最小距離を向上させる。
    \item \textbf{空間結合QLDPC符号}: 特性関数を用いた代数的な定式化により、スタビライザーの可換性を保証しつつ、GRADE-AOアルゴリズム等で短いサイクルを体系的に最適化する。
\end{itemize}

\section*{具体的な構成手順}

提案する量子空間結合（SC）符号の構成は、以下の4つのステップに従う。

\begin{enumerate}
    \item \textbf{マクロ構造（2D-SC HGP）の定義} \\
    Yangらの手法に基づき、2つの古典多項式行列 $A(U,V)$ および $B(U,V)$ を定義する [cite: 601, 755]。これらの行列のテンソル積を用いることで、巨大な親パリティ検査行列 $\hat{H}_X$ および $\hat{H}_Z$ の多項式モデル（特性関数 $F(U,V)$）を構築する [cite: 60, 621]。このマクロな枠組みにより、量子符号としての基本的な構造を担保する。

    \item \textbf{アクティブ領域と潜在領域の分割} \\
    親行列 $\hat{H}_X, \hat{H}_Z$ から、実際に誤り訂正に使用する行を「アクティブ部分（$H_X, H_Z$）」、削除する行を「潜在部分（$\tilde{H}_X, \tilde{H}_Z$）」として分割する [cite: 1596, 1597]。この際、Kasaiの定理に従い、結合長 $L$ に対してアクティブなブロック数 $J$ が $L \ge 4J$ を満たすように設定することで、潜在的な低重み論理演算子の発生を抑制するための数学的マージンを確保する [cite: 1654]。

    \item \textbf{差分集合 $\Delta$ の算出と局所的可換性の適用} \\
    アクティブな行同士が重なり合う多項式の遅延（シフト量）の差分集合 $\Delta$ を算出する [cite: 1642]。$A(U,V)$ と $B(U,V)$ の各要素にアフィン置換行列（APM）を割り当てる際、$\Delta$ の範囲内で重なるブロック同士は可換（内積0）、範囲外で重なるブロック同士は非可換（内積1）となるよう、APMのパラメータ $(a, b)$ を合同式を用いて設計する [cite: 1586, 1692, 1695]。これにより、直交性の制約を局所化し、設計の自由度を向上させる。

    \item \textbf{GRADEアルゴリズムによる最終最適化} \\
    局所的可換条件を満たすAPMの組み合わせ候補群に対し、YangらのGRADE（Gradient-Descent Distributor）アルゴリズムを適用する [cite: 959, 1017, 1023]。勾配降下法を用いて確率分布を最適化した後、AO（Algorithmic Optimization）ステップにより、タナーグラフ上の残存する短いサイクル（長さ6や8など）を最小化するよう、エッジの配置を最終決定する [cite: 1019, 1049]。
\end{enumerate}

\section{マクロ構造：補多項式を用いた親行列の定義}

空間結合量子LDPC符号の設計において、第一段階は、成分行列の具体的な値に依存せず、構造そのものによってパウリ演算子間の可換性を保証する巨視的な枠組みを構築することである。本章では、2変数多項式とハイパーグラフ積（HGP）を用い、遅延の不整合を解決するための「補多項式」を導入することで、符号の母体となる親パリティ検査行列 $\hat{H}_X$ および $\hat{H}_Z$ を定式化する。

\subsection{空間結合の多項式表現とベース行列}
空間結合符号におけるブロックの「ズレ（シフト）」を代数的に扱うため、時間方向の遅延を表す変数 $U$ と、空間方向の遅延を表す変数 $V$ を導入する。量子符号の土台として、以下の2つの古典空間結合符号の設計図となる多項式行列を定義する。

\begin{equation}
    A(U,V) = \sum_{i=0}^{m_1} \sum_{j=0}^{m_2} A_{i,j} U^i V^j \in \mathbb{F}_2^{r_1 \times n_1}[U,V]
\end{equation}
\begin{equation}
    B(U,V) = \sum_{i=0}^{m_1} \sum_{j=0}^{m_2} B_{i,j} U^i V^j \in \mathbb{F}_2^{r_2 \times n_2}[U,V]
\end{equation}

ここで、$m_1$ および $m_2$ はそれぞれ時間方向および空間方向の最大結合長（メモリサイズ）である。係数行列 $A_{i,j}$ と $B_{i,j}$ は、後の段階でアフィン置換行列（APM）が代入されるプレースホルダーとして機能する。

\subsection{直交性の破綻と補多項式の導入}
通常のハイパーグラフ積を空間結合に直接適用し、単なる転置行列 $A(U,V)^T$ や $B(U,V)^T$ を用いて量子パリティ検査行列を構成した場合、シンプレクティック内積の計算において遅延の不整合が発生する。具体的には、$U^i V^j$ と $U^{-i} V^{-j}$ のような正負の遅延が混在し、テンソル積の非対角成分が完全に相殺されず、直交性が破綻する。

この問題を解決するため、多項式 $A(U,V)$ に対する補多項式（Complementary Polynomial） $\bar{A}(U,V)$ を以下のように定義する。

\begin{equation}
    \bar{A}(U,V) = U^{m_1}V^{m_2}A(U^{-1}, V^{-1})
\end{equation}

この変換は、多項式内のすべての遅延を反転（$U \to U^{-1}, V \to V^{-1}$）させた上で、最大遅延 $m_1, m_2$ を乗じて全体を正の遅延空間にシフトし直す操作である。$B(U,V)$ についても同様に $\bar{B}(U,V)$ を定義する。

\subsection{親パリティ検査行列 $\hat{H}_X, \hat{H}_Z$ の構成}
定義した多項式と補多項式を用い、単位行列 $\mathbf{I}$ とのクロネッカー積（$\otimes$）によって、量子スタビライザー符号の親パリティ検査行列 $\hat{H}_X$ および $\hat{H}_Z$ の特性関数を構築する。

\begin{equation}
    \hat{H}_X(U,V) = \begin{bmatrix} \mathbf{I}_{n_2} \otimes A(U,V) & \bar{B}(U,V)^T \otimes \mathbf{I}_{r_1} \end{bmatrix}
\end{equation}
\begin{equation}
    \hat{H}_Z(U,V) = \begin{bmatrix} B(U,V) \otimes \mathbf{I}_{n_1} & \mathbf{I}_{r_2} \otimes \bar{A}(U,V)^T \end{bmatrix}
\end{equation}

ここで、$\mathbf{I}_{n_2}$ はサイズ $n_2 \times n_2$ の単位行列である。このブロック行列表現により、物理量子ビット数は空間結合の各単位ステップあたり $n_1 n_2 + r_1 r_2$ 個となる。

\subsection{巨視的直交性の数学的証明}
構築された親行列が量子符号の必須条件である可換性を満たすことを証明する。直交性は $\hat{H}_X(U,V) \cdot \hat{H}_Z(U^{-1}, V^{-1})^T = \mathbf{0}$ が成立するかで判定される。ブロック行列の乗算を行うと、直交性の鍵となる非対角成分には以下の2つの項の和が現れる。

\begin{equation}
    B(U^{-1}, V^{-1})^T \otimes A(U,V) + \bar{B}(U,V)^T \otimes \bar{A}(U^{-1}, V^{-1})
\end{equation}

補多項式の定義を第二項に代入して展開すると、各多項式に付与されていた遅延シフト $U^{m_1}V^{m_2}$ とその逆数 $U^{-m_1}V^{-m_2}$ が相殺され、$U^0 V^0 = 1$ となって遅延のズレが消滅する。結果として、これら2つの項はテンソル積の構造上まったく同じ多項式行列となり、ガロア体 $\mathbb{F}_2$ 上の演算規則 $1+1=0$ に従って完全に相殺される。

したがって、内部の係数行列 $A_{i,j}, B_{i,j}$ の具体的な値やミクロな配置に依存することなく、マクロなテンソル積と補多項式の構造それ自体によって $\hat{H}_X \hat{H}_Z^T \equiv \mathbf{0}$ が普遍的に保証される。

\section{アクティブ領域と潜在領域の分割による論理演算子の抑圧}

空間結合ハイパーグラフ積（SC-HGP）の巨視的な枠組みにおいて、アフィン置換行列（APM）を導入して剛体サイクルを破壊するためには、成分行列同士の「非可換性」を許容する必要がある。しかし、親行列全体に対して直交性を課すと、APMのパラメータ空間が著しく制限されてしまう。本章では、パリティ検査行列の行を分割し、直交性の制約を局所化するとともに、符号の最小距離を構造的に制限する「潜在的論理演算子」の発生を防ぐメカニズムについて詳述する。

\subsection{親行列からの行分割（パンクチャリング）}
HGP構成法によって構築された巨大な親パリティ検査行列 $\hat{H}_X$ および $\hat{H}_Z$ を、実際にエラーシンドロームの測定に使用する「アクティブ部分（$H_X, H_Z$）」と、符号化や復号の際には使用せず構造から削除される「潜在部分（$\tilde{H}_X, \tilde{H}_Z$）」に水平に分割する。

\begin{equation}
    \hat{H}_X = \begin{bmatrix} H_X \\ \tilde{H}_X \end{bmatrix}, \quad
    \hat{H}_Z = \begin{bmatrix} H_Z \\ \tilde{H}_Z \end{bmatrix}
\end{equation}

この操作は、古典LDPC符号におけるレート調整のための行のパンクチャリング（間引き）に相当するが、量子CSS符号においては、この分割が符号の距離特性に対して極めて重大な影響を及ぼす。

\subsection{潜在的論理演算子化のメカニズムと潜在非直交性}
親行列全体が直交している（$\hat{H}_X \hat{H}_Z^T = \mathbf{0}$）という従来の条件をそのまま維持した状態で上記の分割を行うと、以下の関係式が必然的に成立してしまう。

\begin{equation}
    H_X \tilde{H}_Z^T = \mathbf{0}, \quad H_Z \tilde{H}_X^T = \mathbf{0}
\end{equation}

この式は、削除された潜在行列 $\tilde{H}_Z$ の各行が、アクティブな検査行列 $H_X$ の直交補空間（カーネル）に属することを意味する。つまり、$\tilde{H}_Z$ の低重みな行ベクトルが、そのまま「検出不可能なX論理演算子」へと変貌してしまうのである。これにより、符号の最小距離は削除された行の重み（通常は非常に小さい値）によって構造的に上限が定められてしまう。

この問題を根本から解決するため、本構成法では親行列全体の直交性を意図的に放棄し、以下の「アクティブ直交性」と「潜在非直交性」の双方を同時に要求する。

\begin{equation}
    H_X H_Z^T = \mathbf{0} \quad \text{（アクティブ直交性）}
\end{equation}
\begin{equation}
    H_X \tilde{H}_Z^T \neq \mathbf{0}, \quad H_Z \tilde{H}_X^T \neq \mathbf{0} \quad \text{（潜在非直交性）}
\end{equation}

これにより、潜在部分の行は必ずアクティブな検査行列と干渉し、シンドロームを引き起こすようになるため、論理演算子空間への侵入が物理的に阻止される。これを潜在ベース距離（Latent-based distance）の向上と呼ぶ。

\subsection{結合長とアクティブブロック数の制約}
上記の潜在非直交性を生み出すための数学的な「余白」を確保するためには、結合長に対する制約が生じる。
空間結合におけるブロック行の総数を $L/2$（結合長 $L$ に依存するパラメータ）、実際に使用するアクティブなブロック行数を $J$ と定義する。

親行列全体が直交してしまう（すべてのブロック間でAPMが可換になってしまう）事態を避けるための必要条件として、以下の不等式を満たすよう設計パラメータを選択しなければならない。

\begin{equation}
    L \ge 4J
\end{equation}

この条件を満たすことで、後述するシフト差分の集合 $\Delta$ が全体集合を覆い尽くすことを防ぎ、意図的に「非可換なAPM同士の衝突」を潜在領域に配置することが可能となる。

\subsection{差分集合 $\Delta$ と局所的可換性への帰着}
ブロック巡回構造を持つSC-HGP行列において、行ブロック同士の相互作用は、そのインデックスの差分（シフト量）にのみ依存する。アクティブな行ブロック同士が相互作用し得るシフト差分の集合を $\Delta$ と定義する。

\begin{equation}
    \Delta = \{ (k - i) \pmod{L/2} \mid 0 \le i, k \le J-1 \}
\end{equation}

アクティブ直交性（$H_X H_Z^T = \mathbf{0}$）を満たしつつ、潜在非直交性を達成するための十分条件は、多項式行列に割り当てられるAPM群が以下の規則に従うことである。

\begin{itemize}
    \item \textbf{$\Delta$ に属するシフト差分で重なるAPMペア：} 必ず可換（シンプレクティック内積が0）となるようパラメータ $(a,b)$ を設定する。
    \item \textbf{$\Delta$ に属さないシフト差分で重なるAPMペア：} 少なくとも1つのペアが非可換（シンプレクティック内積が1）となるよう意図的に設計する。
\end{itemize}

このように、直交性の要求を $\Delta$ という局所的な範囲に限定することで、APMの傾き $a$ を柔軟に設定する自由度を獲得し、タナーグラフ上の剛体サイクルを効果的に破壊することが可能となる。

\section{差分集合 $\Delta$ の算出と局所的可換性の適用}

空間結合符号や一般化Hagiwara-Imai符号のようなブロック巡回（Block-circulant）構造を持つパリティ検査行列において、直交性の制約は行列全体ではなく、ブロックインデックスの「差分」に着目することで局所化できる。本章では、アクティブ直交性を満たすための必要十分な相互作用範囲である「差分集合 $\Delta$」を算出し、アフィン置換行列（APM）を用いてその範囲内にのみ可換性を適用する手法を定式化する。

\subsection{ブロック巡回構造と相互作用行列 $\Psi_r$}
親パリティ検査行列 $\hat{H}_X$ および $\hat{H}_Z$ は、サイズ $P \times P$ の置換行列 $F_i, G_i$ を要素とするブロック巡回行列として構成される。ブロック行の総数を $L/2$ としたとき、$\hat{H}_X$ の左半分は $(F_0, F_1, \dots, F_{L/2-1})$ の巡回シフト、右半分は $(G_0, G_1, \dots, G_{L/2-1})$ の巡回シフトから成る。$\hat{H}_Z$ はそれらの転置を用いて対称に構成される。

この構造により、親行列同士の積 $\hat{H}_X (\hat{H}_Z)^T$ の $(i, k)$ ブロック成分は、行インデックス $i$ と列インデックス $k$ の差分 $r := (k - i) \pmod{L/2}$ にのみ依存する。この相互作用行列を $\Psi_r$ と定義すると、以下のように表される。

\begin{equation}
    (\hat{H}_X (\hat{H}_Z)^T)_{i,k} = \sum_{u=0}^{L/2-1} \left( F_u G_{r-u} + G_{r-u} F_u \right) := \Psi_r
\end{equation}

ここで、すべての $u \in [L/2]$ について $F_u$ と $G_{r-u}$ が可換（すなわち $F_u G_{r-u} = G_{r-u} F_u$）であれば、標数2の体（$\mathbb{F}_2$）上の演算規則 $1+1=0$ により $\Psi_r = \mathbf{0}$ となる。

\subsection{差分集合 $\Delta$ の定義とアクティブ直交性}
パリティ検査に実際に使用する「アクティブ行列」 $H_X, H_Z$ は、親行列の先頭から $J$ 個のブロック行（インデックス $0 \le i, k \le J-1$）を抽出したものである。したがって、アクティブ行列同士の積 $H_X H_Z^T$ において発生し得るインデックスの差分は限られている。この相互作用が及ぶ差分の集合を $\Delta$ と定義する。

\begin{equation}
    \Delta := \{ (k - i) \pmod{L/2} \mid 0 \le i, k \le J-1 \} \subset [L/2]
\end{equation}

アクティブ直交性（$H_X H_Z^T = \mathbf{0}$）を満たすためには、すべての $r \in \Delta$ について $\Psi_r = \mathbf{0}$ となれば十分である。すなわち、行列全体を直交させる（すべての $r \in [L/2]$ で $\Psi_r = \mathbf{0}$ とする）必要はなく、$\Delta$ に含まれる差分についてのみ可換性を要求すればよい。

逆に、潜在非直交性（$H_X \tilde{H}_Z^T \neq \mathbf{0}$ など）を達成し、潜在的論理演算子を排除するためには、$\Delta$ に含まれない差分 $r \notin \Delta$ のうち、少なくとも1つについて $\Psi_r \neq \mathbf{0}$ とならなければならない。

\subsection{APMによる局所的可換性と意図的非可換性の設計}
要素となる置換行列 $F_i, G_j$ として、アフィン置換行列（APM）を採用する。サイズ $P$ の素数体（または素数冪）において、各APMは一次関数 $f_i(x) = a_i x + b_i \pmod P$ および $g_j(x) = c_j x + d_j \pmod P$ として表現される。

2つのAPMが可換であるための必要十分条件は、以下の二次合同式で与えられる。

\begin{equation}
    b_i (c_j - 1) \equiv d_j (a_i - 1) \pmod P
\end{equation}

差分 $r$ に寄与するインデックスのペア集合を $\Gamma_r := \{ (u, r-u) \mid u \in [L/2] \}$ と定義する。APMのパラメータ $(a_i, b_i, c_j, d_j)$ の探索においては、以下の2つの条件を同時に満たすように設計を行う。

\begin{enumerate}
    \item \textbf{局所的可換性の適用（$\Delta$ 内部）：} $r \in \Delta$ となるすべてのペア $(i, j) \in \Gamma_r$ について、上記の合同式が成立するようにパラメータを拘束する。
    \item \textbf{意図的非可換性の確保（$\Delta$ 外部）：} 少なくとも1つの $r \notin \Delta$ に対して、合同式が成立しない（非可換となる）ペア $(i, j) \in \Gamma_r$ を意図的に配置し、$\Psi_r \neq \mathbf{0}$ を確保する。
\end{enumerate}

この代数的な設計手法により、剛体サイクルを破壊するための非可換性（傾き $a_i \neq 1$ の自由度）を維持したまま、アクティブな量子スタビライザー条件を完璧に満たすことが可能となる。

\section{パラメータの決定}

空間結合ハイパーグラフ積（SC-HGP）の巨視的な直交性枠組みと、アフィン置換行列（APM）による微視的な非可換性制御を統合した本方式（APM-SC-HGP符号）において、符号の性能を最大化するためのパラメータ設計手順を定式化する。本手法では、代数的な可換性の探索とグラフ理論的なサイクルの最適化を分離・連携させるため、以下の5段階の階層的プロセスを採用する。

\subsection{巨視的ベース構造の決定}
第一段階として、量子符号の土台となる全体の接続関係と符号化率の基本値を決定する。具体的には、元となる2つの古典多項式行列 $A(U,V)$ および $B(U,V)$ のサイズ（$r_1, n_1, r_2, n_2$）と次数分布を設定する。さらに、空間結合のメモリ長（最大遅延 $m_1, m_2$）および2次元の結合長（繰り返し数 $L_1, L_2$）を決定する。これにより、最終的な物理量子ビットの総数（符号長）の枠組みが確定する。

\subsection{アクティブ領域の切り出しと差分集合 $\Delta$ の確定}
第二段階として、潜在的論理演算子の発生を物理的に阻止するための領域分割を行う。結合長 $L_1, L_2$ に対し、実際にシンドローム測定に使用するアクティブなブロック数 $J_1, J_2$ を、制約条件 $L_1 \ge 4J_1$ および $L_2 \ge 4J_2$ を満たすように設定する。この分割により、アクティブなブロック同士が相互作用し得る遅延の差分集合 $\Delta$ が一意に確定し、次段階で可換性を課すべきターゲット範囲が明確となる。

\subsection{代数空間でのAPMパラメータ解の探索}
第三段階では、成分行列として用いるAPMのサイズ $P$ を決定し、配置可能なAPMのパラメータ空間を探索する。部分的な可換および非可換の混在条件を満たす解が存在するよう、$P$ には複数の素因数を持つ合成数を選定することが推奨される。多項式の係数として割り当てるAPMのパラメータ $(a, b)$ について、差分集合 $\Delta$ の内部で重なる場合は可換（合同式が成立）、$\Delta$ の外部で重なる場合は非可換となるような解の集合（解空間）を代数的に探索し抽出する。

\subsection{GRADEアルゴリズムによるトポロジーの最適化}
第四段階は、抽出されたAPMの解空間を用いて、タナーグラフ上の短いサイクル（剛体サイクルおよび柔軟サイクル）を破壊するプロセスである。APMの配置確率を分布として定義し、タナーグラフ上に発生する長さ6や8のサイクルの期待値を目的関数とする。この目的関数に対して勾配降下法（GRADEアルゴリズム）を適用し、短いサイクルが最も生成されにくいAPMの配置確率分布を最適化する。

\subsection{アルゴリズム的最適化（AO）による最終確定}
最終段階として、GRADEアルゴリズムによって得られた最適な確率分布に従い、各遅延位置に具体的なAPMをサンプリングして確定配置する。その後、局所的なエッジの入れ替えを行う貪欲法ベースのアルゴリズム的最適化（AO: Algorithmic Optimization）を適用する。これにより、確率的に残存してしまった微小なサイクルを限界まで排除し、超大内周（Large Girth）を持つAPM-SC-HGP符号のパリティ検査行列を最終的に完成させる。

\section{APM-SC-HGP符号の具体的パラメータセット}

理論的制約（$L \ge 4J$ や $P$ が合成数など）を満たし、実際にプログラムでの検証やシミュレーションに利用できる具体的なパラメータの組み合わせを2パターン提示する。

\subsection{構成案：高性能実証モデル（Girth-8以上を狙う実践的構成）}
実際に低いエラーフロア（FER $10^{-8}$ 以下）と高い最小距離を達成するための、より大規模なパラメータ構成である。

\begin{itemize}
    \item \textbf{古典ベース行列のサイズ:} $r_1=r_2=3, \ n_1=n_2=4$ （行重み4、列重み3のレギュラーLDPCをベースとする）
    \item \textbf{空間結合メモリ:} $m_1=2, \ m_2=2$
    \item \textbf{アクティブブロック数 ($J$):} $J_1=3, \ J_2=3$
    \item \textbf{結合長 ($L$):} $L_1=12, \ L_2=12$ （マージンを十分に確保）
    \item \textbf{APMのサイズ ($P$):} $P=768$ （$2^8 \times 3$）または $P=120$ （$2^3 \times 3 \times 5$）
\end{itemize}

APMサイズ $P$ に $768$ のような複数の素因数を持つ十分大きな合成数を採用することで、差分集合 $\Delta$ （この場合は $\{0, 1, 2, 10, 11\}$ など広範囲に及ぶ）の内部で可換性を保ちつつ、外部で非可換性を爆発させるような解 $(a,b)$ の組み合わせが代数的に確実に存在するようになる。

\subsection{APMパラメータ $(a, b)$ の具体的な選び方}
パラメータ探索アルゴリズムを実装する際、APMの傾き $a$ と切片 $b$ には以下の制約を設けてランダム探索（または全探索）を行う。

APMの一次関数 $y = ax + b \pmod P$ において、逆関数（転置行列）が存在するためには、$a$ と $P$ が互いに素でなければならない。
例えば検証用トイモデル（$P=12$）の場合、$a$ として選べる値は $\{1, 5, 7, 11\}$ のみである。
この中で $a=1$ （従来の巡回行列）の使用を極力避け、$a=5, 7, 11$ と $0 \le b \le 11$ の組み合わせから、以下の合同式を満たすペアと満たさないペアを仕分ける。

$$b_i (c_j - 1) \equiv d_j (a_i - 1) \pmod{12}$$

この条件を満たす部品のセットを抽出し、YangらのGRADEアルゴリズムによって多項式 $A(U,V), B(U,V)$ の適切な遅延位置に配置していくことで、最適な量子パリティ検査行列が完成する。

\section{APM-SC-HGP符号の詳細な構成アルゴリズム}

本章では、巨視的なハイパーグラフ積の枠組みと微視的なアフィン置換行列（APM）の非可換性制御を統合し、エラーフロアのない超大内周量子LDPC符号を構築するための詳細な構成アルゴリズムを記述する。本アルゴリズムは、入力パラメータの初期化から始まり、代数的な解空間の探索、勾配降下法による確率的最適化、そして決定論的な最終調整を経て、アクティブなパリティ検査行列 $H_X$ および $H_Z$ を出力する5つのフェーズで構成される。

\subsection{フェーズ1：マクロ構造の初期化}
アルゴリズムの入力として、古典ベース行列のパラメータ（$r_1, n_1, r_2, n_2$）、空間結合メモリ（$m_1, m_2$）、および結合長（$L_1, L_2$）を与える。
このパラメータに基づき、未定の係数を持つ2変数多項式 $A(U,V)$ および $B(U,V)$ のスケルトン（配置構造）を生成する。続いて、最大遅延 $m_1, m_2$ を用いて補多項式 $\bar{A}(U,V)$ および $\bar{B}(U,V)$ を算出し、単位行列とのクロネッカー積をとることで、全体の自己相殺性を保証する巨大な親パリティ検査行列 $\hat{H}_X, \hat{H}_Z$ の多項式モデルをメモリ上に構築する。

\subsection{フェーズ2：領域分割と差分集合の導出}
潜在的論理演算子を排除するため、親行列からシンドローム測定に用いるアクティブブロック数 $J_1, J_2$ を指定する。このとき、条件 $L_1 \ge 4J_1$ および $L_2 \ge 4J_2$ が満たされていることを検証する。
条件を満たす場合、アクティブ行列 $H_X, H_Z$ を親行列の先頭 $J_1, J_2$ ブロック行として切り出す。さらに、アクティブブロック同士が相互作用し得る遅延のシフト差分集合 $\Delta$ を数式に従って算出し、次フェーズにおける可換性制約のターゲット範囲を確定する。

\subsection{フェーズ3：APMパラメータの代数的探索}
成分行列として用いるAPMのサイズ $P$ （複数の素因数を持つ合成数）を入力とし、多項式の非ゼロ要素に割り当てるためのパラメータ $(a, b)$ の解空間を代数的に探索する。
すべての可能なパラメータペアに対して二次合同式 $b_i (c_j - 1) \equiv d_j (a_i - 1) \pmod P$ の成否を判定し、以下の2条件を同時に満たす「有効なAPM部品の組み合わせセット」を抽出する。
\begin{enumerate}
    \item 差分集合 $\Delta$ の内部で重なる配置においては、合同式が成立し可換となること。
    \item 差分集合 $\Delta$ の外部で重なる配置においては、少なくとも1箇所で合同式が不成立となり非可換となること。
\end{enumerate}
これにより、アクティブ直交性と潜在非直交性を担保する代数的な制約空間が構築される。

\subsection{フェーズ4：GRADEアルゴリズムによる確率的最適化}
フェーズ3で抽出された有効なAPM部品セットを用い、タナーグラフ上の剛体サイクルおよび柔軟サイクルを破壊するためのトポロジー最適化を実行する。
多項式 $A(U,V), B(U,V)$ の各遅延位置に配置されるAPMの選択肢に対して、初期状態として一様分布などの配置確率を割り当てる。タナーグラフ上に形成される短いサイクル（長さ6や8など）の総数の期待値を目的関数とし、この目的関数を最小化するように勾配降下法（GRADEアルゴリズム）を用いて確率分布を反復的に更新する。勾配の大きさが閾値を下回るか、指定された反復回数に到達した時点で更新を終了する。

\subsection{フェーズ5：AOによる最終確定と行列の出力}
GRADEアルゴリズムによって最適化された確率分布に従い、各遅延位置に対して具体的なAPMのパラメータ $(a, b)$ をサンプリングし、ベース行列の要素を確定する。
配置確定後、グラフ上に偶然残存してしまった微小なサイクルを探索し、局所的なエッジの入れ替え（パラメータの微小な変更）を貪欲的に行うAO（Algorithmic Optimization）ステップを実行する。これ以上のサイクル削減が不可能となった時点で最適化を終了し、得られた要素から最終的なバイナリのアクティブパリティ検査行列 $H_X$ および $H_Z$ を展開して出力する。

\section{フェーズ1：マクロ構造の初期化とベース行列の設計}

本章では、量子空間結合符号（APM-SC-HGP符号）を構築する第一段階として、アフィン置換行列（APM）の微視的な値に依存せず、巨視的な接続構造それ自体によってパウリ演算子間の可換性を保証する「マクロ構造の初期化」プロセスについて詳述する。このフェーズでは、最適な古典ベース行列の決定から始まり、補多項式を用いた巨大な親パリティ検査行列の構築までを行う。

\subsection{古典ベース行列 $A, B$ の設計基準}
マクロ構造のシード（種）となる古典ベース行列 $A$ および $B$ は、最終的な量子符号の符号化率やスタビライザーの重み、ノイズ耐性を決定づける最も重要な入力パラメータである。無作為な生成を避け、以下の基準に従って決定する。

\begin{itemize}
    \item \textbf{符号化率とサイズ設定：} ベース行列のサイズ $r_1 \times n_1$ および $r_2 \times n_2$ は、符号の漸近的な符号化率 $R \approx 1 - r_1/n_1 - r_2/n_2$ を決定する。十分な最小距離と高い符号化率を両立させるため、$r=3, n=4$ から $r=3, n=6$ 程度のサイズを初期値として設定することが望ましい。
    \item \textbf{スタビライザー重みの最小化：} 確率伝搬法（BP）による高精度な復号を可能にし、測定エラーを抑制するため、ベース行列の次数分布は疎（Sparse）に保つ。例えば、列重み3、行重み4のレギュラーLDPC行列を採用することで、量子符号におけるパリティ検査の重みを実用的な範囲（重み7程度）に抑えることができる。
    \item \textbf{ノイズモデルへの適応：} 対象とする量子通信路が脱分極チャネルのような対称ノイズであれば $A = B$ とし、超伝導量子ビットのように位相（Z）エラーが支配的な非対称ノイズであれば、$A$ と $B$ の次数分布を非対称に設計することで、特定のパウリ演算子に対する訂正能力を強化する。
\end{itemize}

\subsection{多項式スケルトンの生成と空間結合への展開}
決定した古典ベース行列の非ゼロ要素を、時間方向の最大遅延 $m_1$ および空間方向の最大遅延 $m_2$ で定義される2変数多項式空間 $U^i V^j$ へと展開する。この展開操作においては、要素を単一の遅延位置に集中させず、利用可能な空間 ($0 \le i \le m_1, 0 \le j \le m_2$) に均等かつランダムに分散させる。

これにより、未定の係数（配置指示のみを示すプレースホルダー）を持つ多項式行列 $A(U,V)$ および $B(U,V)$ のスケルトンが生成される。この分散配置が、空間結合特有の閾値向上（Threshold saturation）をもたらす鍵となる。

\subsection{補多項式の導出}
量子符号のX検査とZ検査を直交させる際、単純な転置操作のみでは遅延の正負が混在し、直交性が破綻する。この不整合を解決するため、最大遅延を用いて遅延方向を反転・シフトさせた「補多項式（Complementary Polynomial）」を導出する。行列 $A(U,V)$ に対する補多項式 $\bar{A}(U,V)$ は以下のように定義される。

\begin{equation}
    \bar{A}(U,V) = U^{m_1}V^{m_2}A(U^{-1}, V^{-1})
\end{equation}

この変換は、多項式内のすべての遅延を逆転させた上で、全体を正の遅延空間に引き戻す鏡合わせの操作である。$B(U,V)$ についても同様に $\bar{B}(U,V)$ を算出する。

\subsection{親パリティ検査行列の構築と自己相殺性の担保}
生成した多項式スケルトンと補多項式を用い、単位行列 $\mathbf{I}$ とのクロネッカー積（$\otimes$）によって、全体の自己相殺性を保証する巨大な親パリティ検査行列 $\hat{H}_X$ および $\hat{H}_Z$ の多項式モデルをメモリ上に構築する。

\begin{equation}
    \hat{H}_X(U,V) = \begin{bmatrix} \mathbf{I}_{n_2} \otimes A(U,V) & \bar{B}(U,V)^T \otimes \mathbf{I}_{r_1} \end{bmatrix}
\end{equation}
\begin{equation}
    \hat{H}_Z(U,V) = \begin{bmatrix} B(U,V) \otimes \mathbf{I}_{n_1} & \mathbf{I}_{r_2} \otimes \bar{A}(U,V)^T \end{bmatrix}
\end{equation}

このマクロ構造において、両者の内積 $\hat{H}_X \hat{H}_Z^T$ を計算すると、非対角成分には $B(U^{-1}, V^{-1})^T \otimes A(U,V)$ と $\bar{B}(U,V)^T \otimes \bar{A}(U^{-1}, V^{-1})$ の和が現れる。補多項式の性質により、これら2つの項に含まれる遅延シフトは完全に一致し、ガロア体 $\mathbb{F}_2$ 上の演算規則 $1+1=0$ に従って互いに完全に相殺される。

したがって、後続のフェーズで内部の係数としていかなるアフィン置換行列（APM）を代入しようとも、この構築枠組み自体によって $\hat{H}_X \hat{H}_Z^T \equiv \mathbf{0}$ が数学的に普遍的に保証される。

\section{アフィン置換行列を用いた2次元空間結合HGP符号の完全構成}

本章では、アフィン置換行列（APM）を構成要素として用いた2次元空間結合ハイパーグラフ積（2D-SC HGP）符号のパリティ検査行列の具体的な構成手続きについて詳述する。本構成法は、トポロジーが最適化された「ミクロな部品（APM）」、ズレを伴って空間に広がる「メゾスケールの配線（空間結合）」、そして量子エラー訂正の直交性を担保する「マクロな枠組み（HGP）」という3つの階層的な数理構造の統合によって実現される。

\subsection{ミクロ構造：アフィン置換行列によるベース行列の拡張}

構成の第一歩は、最適化されたアフィン置換行列（APM）を用いて、古典的なベース行列を微視的な配線ブロックへと拡張するプロセスである。サイズ $r \times n$ の古典ベース行列 $A$ の各非ゼロ要素 $(i, j) \in \text{supp}(A)$ に対して、局所最適化によって決定されたパラメータ $(a_{i,j}, b_{i,j})$ を持つサイズ $P \times P$ のAPM $M(\pi_{i,j})$ を割り当てる。このAPMは、以下の一次合同式によって定義される置換 $\pi$ に基づくバイナリ疎行列である。

\begin{equation}
    y \equiv a_{i,j} x + b_{i,j} \pmod P
\end{equation}

ベース行列の要素が $1$ である箇所はこの $P \times P$ のAPMに置き換えられ、$0$ である箇所は $P \times P$ の零行列 $\mathbf{0}_P$ となる。これにより、元のベース行列 $A$ は、サイズ $(rP) \times (nP)$ のブロック行列 $A_{\text{block}}$ へと展開される。

\begin{equation}
    A_{\text{block}} = \begin{bmatrix} M(\pi_{1,1}) & \dots & M(\pi_{1,n}) \\ \vdots & \ddots & \vdots \\ M(\pi_{r,1}) & \dots & M(\pi_{r,n}) \end{bmatrix}
\end{equation}

この段階で、グラフの短い閉路（剛体サイクル）はAPMの傾き $a$ による「ねじれ」によって物理的に破壊され、優れた内周特性を持つ構成パーツが完成する。

\subsection{メゾ構造：テンソル積による空間結合と遅延の展開}

次に、構築された構成パーツを巨大な空間グリッド上に配置し、ブロック間を跨ぐ空間結合（Spatial Coupling）を形成する。結合長 $L_1, L_2$ で定義される2次元グリッド上の「遅延」は、サイズ $L$ の巡回シフト行列 $S_L$ を用いて表現される。

ある特定の遅延 $(u, v)$ を持つ要素のみを抽出した部分行列を $A_{\text{block}}^{(u,v)}$ とすると、空間全体に展開された空間結合行列 $A_{\text{SC}}$ は、すべての遅延パターンの総和として次のように構築される。

\begin{equation}
    A_{\text{SC}} = \sum_{u=0}^{m_1} \sum_{v=0}^{m_2} A_{\text{block}}^{(u,v)} \otimes \left( S_{L_1}^u \otimes S_{L_2}^v \right)
\end{equation}

この操作により、行列 $A_{\text{SC}}$ のサイズは $(r P L_1 L_2) \times (n P L_1 L_2)$ へと拡大する。行列上では、この結合構造は主対角線から斜め方向に向かって伸びる帯状の構造（Diagonal Band）として現れ、量子ビットを保護する堅牢な網目を形成する。

\subsection{マクロ構造：ハイパーグラフ積（HGP）による量子符号の統合}

最終段階として、生成された巨大な空間結合行列 $A_{\text{SC}}$ と $B_{\text{SC}}$ を結合し、量子エラー訂正のスタビライザー条件である直交性を担保するハイパーグラフ積（HGP）を実行する。パリティ検査行列 $\hat{H}_X$ は、単位行列 $\mathbf{I}$ を用いたクロネッカー積によって以下のように構成される。

\begin{equation}
    \hat{H}_X = \begin{bmatrix} \mathbf{I}_{n_2} \otimes A_{\text{SC}} & B_{\text{SC}}^T \otimes \mathbf{I}_{r_1} \end{bmatrix}
\end{equation}

ここで、$\mathbf{I}_{n_2}$ はサイズ $n_2$ の単位行列であり、左側のブロックでは $A_{\text{SC}}$ が $n_2$ 個複製されて対角線上に並ぶ。右側のブロックでは、$B_{\text{SC}}$ の転置行列が $r_1$ 個複製されて配置される。これにより、$\hat{H}_X \hat{H}_Z^T = \mathbf{0}$ の直交条件が厳密に満たされる。

\subsection{本構成法の特性}

本アルゴリズムによる構成は、以下の3つの主要な利点を提供する。

\begin{itemize}
    \item \textbf{エラーフロアの抑制：} APMの傾きにより、空間結合符号特有の致命的な剛体サイクルが排除される。
    \item \textbf{潜在的論理演算子の排除：} 差分集合 $\Delta$ の外部に意図的な非可換領域を設けることで、低重みの論理エラーを破壊できる。
    \item \textbf{閾値の向上：} 2次元的な空間結合による波及効果により、標準的なLDPC符号を上回る高いエラー訂正閾値を達成する。
\end{itemize}

\end{document}